\documentclass[10pt]{article}
\usepackage[T1]{fontenc}
\usepackage[utf8]{inputenc}
\usepackage[french]{babel}
\usepackage{amsmath,amssymb,mathrsfs,mathpazo}
\usepackage{geometry}
\geometry{a4paper,margin=2cm}
\usepackage{graphicx,array,tabularx,tabularray,colortbl,multirow,multicol,booktabs,enumitem,pifont,ulem,fancybox,diagbox,cancel}
\usepackage{pstricks,pstricks-add,pst-plot,pst-tree,pst-node,pst-text,pst-func,pst-math,pst-bspline,pst-3dplot,pst-eucl,pst-solides3d}
\let\pscalculate\undefined
\usepackage{pst-bezier}
\usepackage{tikz,pgfplots,tkz-tab}
\usetikzlibrary{arrows,arrows.meta,patterns,calc,positioning,trees,decorations.markings,shapes.geometric,shapes.misc,shapes.arrows,intersections,decorations.pathmorphing,decorations.pathreplacing,backgrounds,shadows,babel}
\pgfplotsset{compat=1.18}
\usepgfplotslibrary{fillbetween}
\usepackage[np]{numprint}
\usepackage{listings,makecell,pas-tableur,twoopt}
\let\barre\undefined
\usepackage[pstricks]{tablvar}
\definecolor{darkgreen}{RGB}{0,100,0}
\providecommand{\red}{\color{red}}
\providecommand{\blue}{\color{blue}}
\providecommand{\green}{\color{green}}
\providecommand{\black}{\color{black}}
\providecommand{\white}{\color{white}}
\providecommand{\gray}{\color{gray}}
\providecommand{\cyan}{\color{cyan}}
\providecommand{\magenta}{\color{magenta}}
\providecommand{\yellow}{\color{yellow}}
\providecommand{\violet}{\color{violet}}
\providecommand{\orange}{\color{orange}}
\providecommand{\euro}{\texteuro}
\providecommand{\EUR}{\texteuro}
\providecommand{\ieme}{\textsuperscript{e}}
\providecommand{\degres}{\ensuremath{^\circ}}
\providecommand{\degree}{\ensuremath{^\circ}}
\providecommand{\vv}[1]{\overrightarrow{#1}}
\providecommand{\Coord}[1]{#1}
\providecommand{\decosix}{}
\providecommand{\starredbullet}{\ensuremath{\bullet}}
\pagestyle{empty}
\setlength{\parindent}{0pt}
\usepackage{adjustbox}\DeclareUnicodeCharacter{2236}{:}\DeclareUnicodeCharacter{2190}{\ensuremath{\leftarrow}}
\providecommand{\R}{\mathbb{R}}
\providecommand{\N}{\mathbb{N}}
\providecommand{\Z}{\mathbb{Z}}
\providecommand{\Q}{\mathbb{Q}}
\providecommand{\D}{\mathbb{D}}
\providecommand{\C}{\mathbb{C}}
\providecommand{\e}{{\mathrm{e}}}
\providecommand{\vect}[1]{\overrightarrow{#1}}
\providecommand{\vectt}[1]{\overrightarrow{\mathrm{#1}}}
\providecommand{\vvt}[1]{\overrightarrow{\mathrm{#1}}}
\providecommand{\np}[2][]{{#2~#1}}
\providecommand{\vv}[1]{{\overrightarrow{#1}}}
\renewcommand{\d}{{\mathrm{d}}}
\providecommand{\degres}{{^\circ}}
\providecommand{\degre}{{^\circ}}
\providecommand{\blue}{\color{blue}}
\providecommand{\red}{\color{red}}
\providecommand{\enskip}{\quad}
\providecommand{\llbracket}{\lbrack\!\lbrack}
\providecommand{\rrbracket}{\rbrack\!\rbrack}
\providecommand{\euro}{{\text{€}}}
\providecommand{\no}{{\text{n°}}}
\providecommand{\interval}[2]{{[#1;#2]}}
\providecommand{\intervFO}[2]{{[#1;#2[}}
\providecommand{\dotfill}{\ldots}

\providecommand{\ds}{\displaystyle}
\providecommand{\textding}[1]{#1}
\providecommand{\Oij}{\(\mathrm{O};\vec{\imath},\vec{\jmath}\)}
\providecommand{\Oijk}{\(\mathrm{O};\vec{\imath},\vec{\jmath},\vec{k}\)}
\providecommand{\Ouv}{\(\mathrm{O};\vec{u},\vec{v}\)}
\providecommand{\intervFF}[2]{\ensuremath{[#1;#2]}}
\renewcommand{\euro}{\ifmmode\text{\texteuro}\else\texteuro\fi}
\DeclareUnicodeCharacter{00B0}{\ifmmode^{\circ}\else\textdegree\fi}
\DeclareUnicodeCharacter{20AC}{\ifmmode\text{\texteuro}\else\texteuro\fi}
% Énoncés : Éducation nationale. Transcription et corrigés : APMEP, Denis Vergès et contributeurs. Découpage et classement : CoopMaths, Sébastien Lozano, Éric Elter et contributeurs. Adaptation de la présentation : Maths Explorer.
% CC BY-SA : https://coopmaths.fr/www/about/
% https://coopmaths.fr/alea/static/bac/2021/tex/bac_2021_06_etrangers_7.tex
\begin{document}
\noindent\textit{APMEP / CoopMaths. CC BY-SA. Présentation adaptée par Maths Explorer.}\par\medskip
\textbf{Commun à tous les candidats}

\medskip

Au 1\up{er} janvier 2020, la centrale solaire de Big Sun possédait \np{10560} panneaux solaires.

On observe, chaque année, que 2\,\% des panneaux se sont détériorés et nécessitent d'être retirés tandis que 250 nouveaux panneaux solaires sont installés.

\bigskip

\textbf{Partie A - Modélisation à l'aide d'une suite}

\medskip

On modélise l'évolution du nombre de panneaux solaires par la suite $\left(u_n\right)$ définie par $u_0 = \np{10560}$ et, pour tout entier naturel $n$,\, $u_{n+1} = 0,98u_n +250$, où $u_n$  est le nombre de panneaux solaires au 1\up{er} janvier de l'année $2020 +n$.

\medskip

\begin{enumerate}
\item 
	\begin{enumerate}
		\item  Expliquer en quoi cette modélisation correspond à la situation étudiée.
		\item  On souhaite savoir au bout de combien d'années le nombre de panneaux solaires sera strictement supérieur à \np{12000}.
		
À l'aide de la calculatrice, donner la réponse à ce
problème.
		\item  Recopier et compléter le programme en Python ci-dessous de sorte que la valeur
cherchée à la question précédente soit stockée dans la variable n à l'issue de l'exécution de ce dernier.
\begin{center}
\begin{tabular}{|l|}\hline
u $= \np{10560}$\\
n$ =0$\\
while  \ldots \ldots :\\
\qquad u $= \ldots \ldots$ \\
\qquad n $= \ldots \ldots$\\ \hline
\end{tabular}
\end{center}

	\end{enumerate}
\item Démontrer par récurrence que, pour tout entier naturel $n$, on a $u_n \leqslant \np{12500}$.
\item Démontrer que la suite $\left(u_n\right)$ est croissante.
\item En déduire que la suite $\left(u_n\right)$ converge. Il n'est pas demandé, ici, de calculer sa limite.
\item On définit la suite $\left(v_n\right)$  par $v_n = u_n - \np{12500}$, pour tout entier naturel $n$.
	\begin{enumerate}
		\item Démontrer que la suite $\left(v_n\right)$ est une suite géométrique de raison $0,98$ dont on précisera le premier terme.
		\item Exprimer, pour tout entier naturel $n$,\, $v_n$ en fonction de $n$.
		\item En déduire, pour tout entier naturel $n$,\, $u_n$ en fonction de $n$.
		\item Déterminer la limite de la suite $\left(u_n\right)$. 
		
Interpréter ce résultat dans le contexte du modèle.
	\end{enumerate}
\end{enumerate}

\bigskip

\textbf{Partie B - Modélisation à l'aide d'une fonction}

\medskip

Une modélisation plus précise a permis d'estimer le nombre de panneaux solaires de la centrale à l'aide de la fonction $f$ définie pour tout $x \in [0~;~ +\infty[$ par 

\[f(x) = \np{12500} - 500 \e^{-0,02x+1,4},\]

 où $x$ représente le nombre d'années écoulées depuis le 1\up{er} janvier 2020.

\medskip

\begin{enumerate}
\item Étudier le sens de variation de la fonction $f$.
\item Déterminer la limite de la fonction $f$ en $+\infty$.
\item En utilisant ce modèle, déterminer au bout de combien d'années le nombre de panneaux
solaires dépassera \np{12000}.
\end{enumerate}

\bigskip

\end{document}
